\section{Related Work}
\label{sec:relworks}

Different VLMs and Tiny LLM architectures have emerged that enable deployment and applications of multimodal AI on resource-constrained devices. Recent developments in small VLMs, such as H2OVL-Mississippi (0.8B parameters)~\citep{galib-2024-h2ovlmississippivisionlanguagemodels}, TinyGPT-V~\citep{tinygpt-v-2024}, and MiniCPM-V~\citep{minicpm-v-2024}, demonstrate that compact multimodal models can achieve competitive performance while maintaining efficient deployment characteristics. Similarly, Tiny LLMs, such as MobileLLM~\citep{mobilellm-2024} and TinyLLM~\citep{tinyllm-edge-2024}, have shown that sub-billion parameter models can be quantized and optimized for their deployment on edge devices. These highlight the feasibility of on-device multimodal processing, with models providing meaningful performance while addressing security, latency, and connectivity constraints.

Furthermore, memory-augmented neural networks and language models inspired by cognitive thinking, such as humans, have also garnered significant attention for their ability to store and retrieve contextual information related to specific things across certain short periods of time. 
Memory-augmented neural networks (MANNs)~\citep{graves-dynamicexternalmemory-2016}, use decoupled key-value structures to store and retrieve contextual information. Recent works, such as EGO~\citep{ego-episodic-2024} and selective episodic memory strategies~\citep{episodic-memory-neural-2022}, have extended these ideas for flexible knowledge transfer and context-based memory access. However, these models face limitations in combining memory systems with low-bit quantized multimodal encoders, often sacrificing either memory capacity or model precision.
% And recent works on similar episodic memory includes the Episodic Generalization and Optimization (EGO) framework, which integrates the episodic memory module, with semantic pathways to enable rapid learning and flexible knowledge transfer across tasks~\citep{ego-episodic-2024}. And Memory-Augmented Neural Networks (MANNs)~\citep{graves-dynamicexternalmemory-2016} have demonstrated improved performance through decoupled key-value structures that preserve explainability of historical knowledge extraction~\citep{decoupled-memory-2024}, while research on selective episodic memory encoding and retrieval shows that neural networks can learn human-like strategies for determining when to store and access memories based on environmental uncertainty and contextual relevance~\citep{episodic-memory-neural-2022}.

BitMar overcomes these challenges by integrating 1.58-bit quantization across text and vision encoders, alongside a cross-modal memory retrieval system. The design enables BitMar to store and retrieve both textual and visual context, improving memory interactions and enhancing multimodal generation tasks, all while maintaining computational efficiency for edge deployment.
% Despite these advances in tiny multi-modal models, and human-inspired memory systems, no existing work has successfully combined ultra-low-bit quantized multi-modal encoding with episodic memory mechanisms in a unified architecture optimized for edge deployment. The current models either sacrifice their memory capabilities for efficiency, or they sacrifice their precision, to produce skewed results. BitMar, thus addresses this gap by integrating an 1.58-bit quantization across both text and vision encoders with a human-like episodic memory system.
